\documentclass[11pt,letterpaper]{article}
% Shared layout. Compile each document from this directory.
\usepackage[margin=0.8in,headheight=14pt,headsep=16pt,footskip=22pt]{geometry}
\usepackage[T1]{fontenc}
\usepackage{lmodern,amsmath,amssymb,booktabs,enumitem,titlesec,fancyhdr,microtype,xcolor,hyperref}
\definecolor{inkblue}{RGB}{25,65,100}
\hypersetup{colorlinks=true,linkcolor=inkblue,urlcolor=inkblue,filecolor=inkblue,pdfauthor={AI-assisted course study notes}}
\setlength{\parindent}{0pt}
\setlength{\parskip}{5pt}
\setlength{\emergencystretch}{1.5em}
\setlength{\abovedisplayskip}{7pt}
\setlength{\belowdisplayskip}{7pt}
\titleformat{\section}{\normalsize\bfseries\color{inkblue}}{}{0pt}{}
\titlespacing*{\section}{0pt}{9pt}{3pt}
\setlist[itemize]{leftmargin=15pt,itemsep=2pt,topsep=3pt,parsep=0pt}
\setlist[enumerate]{leftmargin=17pt,itemsep=3pt,topsep=3pt,parsep=0pt}
\pagestyle{fancy}
\fancyhf{}
\fancyhead[L]{\small ECE 750T\enspace /\enspace Responsible AI}
\fancyhead[R]{\small Fall 2026\enspace /\enspace Study notes}
\fancyfoot[R]{\small\thepage}
\renewcommand{\headrulewidth}{0.3pt}
\newcommand{\lesson}[2]{%
  \hypersetup{pdftitle={ECE 750T: #1 -- #2}}%
  {\LARGE\bfseries #1}\par\vspace{3pt}
  {\large #2}\par\vspace{3pt}
  {\footnotesize AI-assisted reconstruction from course slides, not a lecture transcript.
  Discussion reasoning is inferred; classroom responses were not supplied.}\par
  \vspace{3pt}\hrule\vspace{5pt}}
\newcommand{\sources}[1]{\par\vspace{3pt}{\footnotesize\textbf{Slide trail.} #1
  Page numbers refer to PDF pages, not the slides' continuous printed numbering.}\par}
\newcommand{\Prob}{\mathbb{P}}
\newcommand{\E}{\mathbb{E}}
\newcommand{\Yhat}{\widehat{Y}}
\newcommand{\indep}{\mathrel{\perp\!\!\!\perp}}

\begin{document}
\lesson{Midterm study guide}{Confirmed scope, reasoned priorities, and worked checks}

\section*{What the supplied documents actually confirm}
L01 p.~17 specifies Weeks 1--5, conceptual questions and worked problems, with
the exam on October 23. The midterm is worth \textbf{10\% of the course grade}.
The cover specifies closed book with a provided reference sheet and 28 marks:
Problem 1 has eight parts worth eight marks; Problem 2 has four parts worth four
marks; Problem 3 requires \emph{two of three} short answers worth 16 marks total.
There are no part marks for Problems 1--2 and no penalty for incorrect guesses.
Problem 3 permits part marks and explicitly rewards clarity and concision.
If all three choices are answered, only 3(a) and 3(b) are marked.

The cover does not identify topics or the response format for Problems 1--2.
No question paper was supplied. These are study priorities, not predictions of questions.

\section*{Where to spend the most effort}
\begin{enumerate}
\item \textbf{Fairness calculations and interpretation (L02).} The reference sheet
explicitly supplies confusion-matrix rates, group criteria, calibration, and the
prevalence identity. Practise denominators, conditioning directions, incompatible
criteria, and the assumptions and exceptions. Also review individual fairness,
earthmover distance, and A1's smooth-versus-hard parity distinction.
\item \textbf{Privacy guarantees and counterexamples (L05).} The sheet explicitly
supplies $k$-anonymity, add/remove adjacency, pure/approximate DP, composition,
and post-processing. Practise applying them. Review sensitivity, Laplace noise,
privacy units, trust boundaries, and why anonymization or federated learning
alone does not establish privacy.
\item \textbf{Concise scenario analysis (L01, L03, L04, and L02's readings).}
Short answers are $16/28\approx57\%$ of the marks. Prepare to diagnose a target
or proxy, evaluation flaw, abstraction trap, disclosure risk, or missing audit
evidence, and propose a response whose effectiveness could be checked.
\end{enumerate}

\textbf{The sheet's omissions do not exclude other Weeks 1--5 material.}
Review agency, power, intersectional evaluation, benchmark exposure, validity,
documentation, reward hacking, estimation SE, and the finite-class bound.
News dates, biographies, software defaults, exact rankings, and advanced methods
absent from the decks are lower priorities. Topics only announced as future
material have weaker support than concepts developed in L01--L05.

\section*{Worked fairness check: equal errors, different access and PPV}
\begin{center}
\small
\begin{tabular}{@{}lrrrrrrrr@{}}
\toprule
Group & TP & FN & FP & TN & Prevalence & TPR & FPR & PPV \\
\midrule
$a$ & 40 & 10 & 10 & 40 & .50 & .80 & .20 & .80 \\
$b$ & 16 & 4 & 16 & 64 & .20 & .80 & .20 & .50 \\
\bottomrule
\end{tabular}
\end{center}
Both groups have 100 people. Award rates are $.50$ and $.32$, so the DP gap is
$.18$. Equalized odds and equal opportunity hold; demographic parity and predictive
parity fail. The difference follows from $(.8-.2)(.5-.2)=.18$.
This establishes statistical comparisons relative to the recorded $Y$, not that
the labels are valid or that the award policy is justified.

\newpage
\section*{Worked privacy checks: apply the quantifiers}
\textbf{Anonymity.} Suppose the quasi-identifier class sizes are $3,2,4$.
The largest $k$ satisfied is two, the minimum class size. Adding a new
quasi-identifier can split classes and reduce $k$. If all diagnoses within a
matching class are identical, knowing membership reveals the diagnosis despite
row ambiguity.

\textbf{An event can refute a DP claim.} For neighboring $D,D'$ and an event
$S$, suppose probabilities are $.30$ and $.10$. Pure DP requires
$.30\le e^\epsilon(.10)$, hence $\epsilon\ge\log 3$ for this ordered pair
and event. At $\epsilon=\log2$, approximate DP requires $\delta\ge.10$.
Satisfying this one constraint does not certify the mechanism: all neighboring
pairs, both orders, and all events must work. A positive-probability event
opposed to a zero-probability event rules out any finite pure-DP parameter.

\textbf{Noise and repeated releases.} A sensitivity-one count with
$\epsilon=.5$ uses Laplace scale $b=2$, has mean absolute error two, and
a 95\% absolute-error bound $2\log20\approx5.99$.
Three fresh releases each $(.5,10^{-6})$-DP give a basic transcript bound
$(1.5,3\!\times\!10^{-6})$-DP. Publishing a rounded version of those same
outputs adds no private-data access; recomputing fresh noisy counts does.

\section*{A1-style reasoning: do not overgeneralize a slogan}
Independence can disappear after conditioning. For independent fair bits $X,R$
and $Y=X\mathbin{\mathrm{xor}}R$, conditioning on $Y=0$ forces $X=R$:
$\Prob(X=0,R=0\mid Y=0)=1/2$, whereas the product of conditional marginals
is $1/4$. Thus $X\indep R$ but $X\not\indep R\mid Y$.

The binary-outcome conflict is also assumption-dependent. A1 explicitly asks
about a three-valued target. For $Y\in\{0,1,2\}$, let group-conditional
probabilities be $(.4,.1,.5)$ and $(.1,.4,.5)$, and let
$R=\mathbf1\{Y=2\}$. Then $A$ and $Y$ are dependent, but $R$ has rate .5
in both groups and $R\indep A\mid Y$. Separation and independence can coexist.
Likewise, a constant predictor can satisfy both without being useful.
For A1's group-prediction result, set $q_a=\Prob(g(X)=1\mid A=a)$.
Choose $h=g$ if $q_1\ge q_0$ and $h=1-g$ otherwise; then
$\mathrm{RWAcc}_A(h)=\tfrac12+\tfrac12|q_1-q_0|$.

\section*{A concise answer pattern for the conceptual questions}
State a claim, give the mechanism and relevant assumption, connect it to a
consequence, and specify evidence or a justified response. For example:

\emph{``Using spending as health need is a validity concern because access to
care affects spending independently of need. A model can predict costs accurately
while missing unmet need in a poorly served group. Evaluate clinical-need measures
and subgroup allocation consequences before treating this score as a need measure.''}

Practise applying that pattern to these slide-based prompts:
\begin{itemize}
\item \textbf{A new-hospital claim from a random image split:} use a held-out site;
patient overlap is a separate leakage concern. Explain which population the score estimates.
\item \textbf{A developer-funded embedded audit:} access may improve, while
publication rights, independence, and authority to require action still need evidence.
\item \textbf{A model card proposed for a new population:} identify the relevant
evaluation gap and evidence needed; do not invent subgroup failures.
\item \textbf{An autonomous scholarship agent:} examine its authority, objective,
feedback, affected people, appeals, and who can change the policy or stop it.
\end{itemize}

\sources{Confirmed facts: \href{midterm-cover.pdf}{midterm-cover.pdf}, p.~1;
\href{reference-sheet.pdf}{reference-sheet.pdf}, p.~1;
\href{../lectures/L01.pdf}{L01.pdf}, pp.~13, 17. Topic priorities draw on all five decks
and A1/A2 handouts. Worked numerical examples here are study constructions,
not released exam questions. Build from the repository root with Tectonic:
\texttt{sh lectures/summaries/build.sh}. Each PDF has a matching \texttt{.tex} source.}
\end{document}
